\section{Discussion}
The clearest result is the improvement obtained by accounting for production conditions. A regularized linear model reduced error against the machine-level intensity reference on all three machines in a later industrial period. This gives the baseline a practical role in comparing completed operations with different masses and product characteristics. The sizeable Pellet~1 underestimation nevertheless shows that a good global result does not remove the need to examine individual machines. The available records do not resolve whether that bias arises from omitted operating conditions, product allocation, changing consumption or measurement boundaries. A baseline fitted to historical operation also reflects any inefficiencies already present in that operation; it does not estimate minimum achievable energy use.

The temporal comparison suggests that, in this dataset, current production characteristics carry much of the usable predictive information. TFT and the current-input Ridge control were close, whereas the full-window Ridge and N-HiTS had larger errors. This pattern is consistent with limited additional value from the tested histories, but it does not identify a causal explanation for an architecture's performance. Each event aggregates a different mass over a different duration, and the event index does not represent a uniform time grid. Moreover, the strict history reset removed 27\% of the original validation targets as warm-up, narrowing the evaluated population. Longer histories, elapsed-time representations and broader architecture selection could change the comparison, but would require their own development protocol. The present fixed configurations do not establish a general ordering of forecasting model families.

The foundation-model diagnostics sharpen this point. Covariate adaptation and normalization influence the final industrial prediction alongside the pretrained weights. TimesFM XReg extrapolated severely when its local mass history was nearly constant; Chronos-2 with covariates showed sensitivity to device-dependent normalization of constant channels (Appendix~\ref{sec:computing}). These observations identify concrete issues to investigate before operational use. They also explain why merely making a model accept the available covariates is insufficient evidence that the complete prediction pipeline is stable for this production setting.

The small exploratory gains from product-dependent mass effects and residual correction offer a more focused next step. Product subfamily can partly distinguish operating demands that a single mass coefficient averages together. Recent residuals may capture persistent departures from the historical relation, provided that updates use only completed outcomes and are sufficiently stabilized. Neither interpretation establishes a causal product effect or equipment change. The favorable results support prespecifying these two candidates for a new period, with separate assessment of machine-level bias, rather than replacing the frozen model on the basis of additional comparisons with an already examined test.

For prospective forecasting, the main unresolved issue is when the inputs become known. The consolidated workbook contains actual dosed mass and batch totals, without plan versions or publication timestamps. Those fields are suitable for retrospective production adjustment, but may only become complete after processing has started. Overlapping records further complicate the relation between a source row and a physical production boundary. Confirming meter scope, stable order identifiers and outcome publication times would therefore add more value to the next evaluation than further refinement of an unverified timestamp proxy. Real time-resolved observations would also be needed to test multiple clock-time horizons, machine inactivity or within-operation transitions. The current one-record task cannot resolve those questions.

Uncertainty should be judged by what the intervals enable in practice. Here, updating recent calibration increased coverage largely through wider bounds and gave little improvement in interval score. The trade-off limits the value of treating every out-of-interval record as an operational signal. A new consecutive industrial period should retain all eligible operations starting within it and follow long operations to completion, with the baseline, adjustment rules and calibration policy fixed beforehand. Such a protocol would provide the evidence needed to integrate the baseline into a digital twin \citep{tao2019}. The interface could retain the input version and model version for each issued estimate, then attach the measured outcome when it becomes available. This would preserve a usable distinction between an earlier forecast and a later production-adjusted assessment, while supporting plant review of unusually high consumption.
